% ============================================================
% ANEXO: Manual de Fabricación, Instalación y Operativa
% Sin \appendix aquí — va en main.tex
% ============================================================

% ============================================================
\chapter{Fabricación del Nodo Radar}
\label{anexo:fabricacion}
% ============================================================

En esta parte del documento se detallan los pasos necesarios para construir un nodo radar periférico, grabar su firmware y dejarlo listo para conectarse al sistema.

\section{Requisitos Previos}

Para la fabricación de cada nodo se requiere:

\begin{itemize}
    \item Piezas estructurales impresas en 3D (chasis inferior, placa intermedia, placa superior, soportes y cabezal rotatorio).
    \item Motor paso a paso NEMA~17.
    \item Driver A4988.
    \item Sensor ultrasónico HC-SR04.
    \item Microcontrolador ESP32-WROOM-32D.
    \item Módulo regulador DC-DC LM2596.
    \item Reed switch magnético e imán permanente.
    \item Fuente de alimentación centralizada de 24~V.
    \item Protoboard, cables dupont y resistencias de 10~k$\Omega$.
    \item Ordenador con Visual Studio Code y extensión PlatformIO instalada.
\end{itemize}

\section{Montaje de la Estructura Mecánica}

El ensamblaje de las piezas debe realizarse en el siguiente orden:

\begin{enumerate}
    \item Colocar la placa inferior e insertar los pilares de separación en las esquinas.
    \item Fijar el motor NEMA~17 en su soporte sobre la placa intermedia.
    \item Instalar el soporte del reed switch adyacente al motor.
    \item Encajar a presión los soportes del ESP32 y la protoboard en la placa inferior.
    \item Fijar el módulo LM2596 y ajustar su potenciómetro hasta obtener exactamente $5{,}0$~V de salida, verificado con multímetro.
    \item Acoplar el cabezal rotatorio al eje del motor mediante encaje a presión.
    \item Insertar el sensor HC-SR04 en el alojamiento del cabezal.
    \item Fijar el imán en el cabezal, en el lado opuesto al sensor.
    \item Cerrar el chasis con la placa superior, organizando el cableado para que no interfiera con el giro del cabezal.
\end{enumerate}

Las Figuras~\ref{fig:nodo_completo} y~\ref{fig:nodo_montaje} muestran el nodo completamente ensamblado desde dos ángulos complementarios.

\begin{figure}[H]
    \centering
    \includegraphics[width=0.55\textwidth]{fig_anexo_nodo_completo}
    \caption{Vista frontal del nodo radar ensamblado.}
    \label{fig:nodo_completo}
\end{figure}

\begin{figure}[H]
    \centering
    \includegraphics[width=0.55\textwidth]{fig_anexo_nodo_montaje}
    \caption{Vista lateral del nodo radar ensamblado.}
    \label{fig:nodo_montaje}
\end{figure}

\section{Conexiones Electrónicas}

El diagrama completo de conexiones se encuentra en la Figura~\ref{fig:diagrama_conexiones} del Capítulo~\ref{chap:disenyo}. La asignación de pines es la siguiente:

\begin{itemize}[topsep=4pt, partopsep=0pt, itemsep=2pt, parsep=0pt]
    \item \textbf{HC-SR04:} VCC $\rightarrow$ 5\,V;\; TRIG $\rightarrow$ D14;\; ECHO $\rightarrow$ D27;\; GND $\rightarrow$ GND.
    \item \textbf{Reed switch:} Lado 1 $\rightarrow$ D26 (pull-up interno);\; Lado 2 $\rightarrow$ GND.
    \item \textbf{A4988 (lógica):} STEP $\rightarrow$ D19;\; DIR $\rightarrow$ D18;\; MS1, MS2, MS3 $\rightarrow$ 3V3;\; VDD $\rightarrow$ 3V3;\; EN $\rightarrow$ GND.
    \item \textbf{A4988 (potencia):} VMOT $\rightarrow$ Bus 24\,V;\; GND $\rightarrow$ GND.
    \item \textbf{LM2596:} VIN $\rightarrow$ Bus 24\,V;\; VOUT $\rightarrow$ 5\,V;\; GND $\rightarrow$ GND.
    \item \textbf{ESP32:} VIN $\rightarrow$ 5\,V (LM2596);\; GND $\rightarrow$ GND.
\end{itemize}

\section{Grabación del Firmware}

\begin{enumerate}
    \item Abrir la carpeta del proyecto en Visual Studio Code. PlatformIO detectará el archivo \texttt{platformio.ini} automáticamente.

    \item Conectar el ESP32 al ordenador mediante cable USB.

    \item Editar \texttt{src/config.h} y localizar el array \texttt{KNOWN\_NETWORKS}. El proyecto incluye dos redes preconfiguradas. Para añadir nuevas redes, agregar entradas adicionales en el mismo formato:

\begin{verbatim}
{"SSID_red", "contraseña_red"},
\end{verbatim}

    \item Pulsar \textbf{Build} en la barra de PlatformIO para compilar y verificar que no hay errores.

    \item Pulsar \textbf{Upload} para grabar el firmware en el ESP32. Tras la grabación, el nodo realizará el homing y quedará a la espera del servidor.
\end{enumerate}

% ============================================================
\chapter{Instalación y Configuración del Sistema}
\label{anexo:instalacion}
% ============================================================

En esta parte se describe cómo instalar el entorno software, desplegar los nodos físicamente y configurar el sistema antes de la primera sesión.

\section{Requisitos Previos}

\begin{itemize}
    \item Ordenador con Python 3.8 o superior.
    \item Tres nodos radar fabricados y con firmware grabado.
    \item Red WiFi local accesible desde el ordenador y los nodos.
    \item Cinta métrica para la medición de coordenadas físicas.
\end{itemize}

\section{Instalación de Dependencias}

Las únicas dependencias no estándar del sistema son las siguientes. Instalarlas con:

\begin{verbatim}
pip install zeroconf ifaddr
\end{verbatim}

La interfaz gráfica utiliza exclusivamente librerías incluidas en la distribución estándar de Python y no requiere instalación adicional.

\section{Arranque del Sistema}

Ejecutar el siguiente comando desde la carpeta \texttt{tools/} del proyecto:

\begin{verbatim}
python radar_visualizer.py
\end{verbatim}

Se abrirá el panel gráfico en Modo Diseño. El servidor no arranca hasta que se pulse \textbf{GUARDAR Y EJECUTAR}.

\section{Ubicación Física de los Nodos}

Colocar los tres nodos sobre la superficie de vigilancia con el sensor HC-SR04 orientado hacia el interior del área, tal como muestra la Figura~\ref{fig:anexo_despliegue_fisico}. Medir con una cinta métrica las distancias entre nodos para obtener sus coordenadas en centímetros.

\begin{figure}[H]
    \centering
    \includegraphics[width=0.75\textwidth]{fig_anexo_despliegue_fisico}
    \caption{Despliegue físico de los tres nodos radar en disposición triangular sobre la superficie de vigilancia.}
    \label{fig:anexo_despliegue_fisico}
\end{figure}

\section{Configuración de Coordenadas y Zonas de Alerta}

Las coordenadas, ángulos y zonas de alerta se configuran en el Modo Diseño, como se muestra en la Figura~\ref{fig:gui_parametros}. Los nodos se reposicionan arrastrando su icono sobre el canvas; las zonas DEFCON se ajustan con los controles deslizantes del panel lateral. Al pulsar \textbf{GUARDAR Y EJECUTAR} toda la configuración se guarda en disco y el servidor arranca.

% ============================================================
\chapter{Guía de Operativa}
\label{anexo:operativa}
% ============================================================

La interfaz gráfica se describe en detalle en la Sección~5.4 de esta memoria. A continuación se recogen los aspectos operativos esenciales.

\section{Modo Diseño}

Permite configurar la geometría del despliegue y los umbrales de alerta antes de iniciar la vigilancia. Una vez listos todos los parámetros, pulsar \textbf{GUARDAR Y EJECUTAR} para arrancar el servidor y pasar al Modo Telemetría.

\section{Modo Telemetría}

Muestra las detecciones en tiempo real. Cada eco aparece como un punto de color según la zona de alerta alcanzada (rojo DEFCON~1, naranja DEFCON~2, amarillo DEFCON~3, azul claro valla). Cuando dos nodos detectan el mismo objeto en zona de solapamiento, aparece un círculo amarillo con las coordenadas de la posición cooperativa estimada.

El panel lateral muestra el estado de cada nodo. Si un nodo se desconecta, su tarjeta indica \textit{DESCONECTADO / SIN DATOS} hasta la reconexión automática.

Para finalizar la sesión, pulsar \textbf{DETENER SERVIDOR}. El sistema guarda el estado angular de todos los nodos y regresa al Modo Diseño.

\section{Reanudación de Sesión}

Al iniciar el visualizador, si existe un estado de sesión anterior guardado en disco, el sistema ofrece la opción de reanudarlo, recuperando posiciones y ángulos de todos los nodos, o comenzar una sesión nueva.
